% GENERATED FILE. Edit canonical lessons, then run npm run generate:books.
% canonical-source-hash: fnv1a64:f66caae4be79bb79
% canonical-lessons: TA-C217-manamakal, TA-C217-manamakan, TA-C217-niccayatarttam, TA-C217-viruntu, TA-C217-vila

\chapter{Bride, Bridegroom, Engagement, Feast, Festival}
\label{ch:ta-bride-bridegroom-engagement-feast-festival}
\hlchaptermodality{217}

\begin{chapteropening}
\textbf{By the end of this chapter:} \emph{I can say a bride, a bridegroom, an engagement, a feast and a festival.}

Complete the last lesson of chapter 217: I can say a bride, a bridegroom, an engagement, a feast and a festival.
\end{chapteropening}

\section[\texorpdfstring{\ta{மணமகள்} (maṇamakaḷ)}{மணமகள் (maṇamakaḷ)}]{\ta{மணமகள்} (maṇamakaḷ) — a bride}

\label{lesson:TA-C217-manamakal}



\begin{quote}
Before the new one: say the Tamil for a ticket; a slip of paper, then the Tamil for a map.
\end{quote}

\subsection*{You'll want to know: \ta{மணமகள்}}
\textbf{\ta{மணமகள்}} — \emph{maṇamakaḷ} — \textquotedblleft{}a bride\textquotedblright{}.

\begin{grammarlens}[title={What you've built}]
One more word for this chapter.
\end{grammarlens}

\subsection*{Your turn}
\begin{itemize}
\raggedright
  \item \emph{Say it:} \emph{maṇamakaḷ}
  \item \emph{Say it:} \emph{maṇamakaḷ}, once more
  \item \emph{Say it:} say \emph{maṇamakaḷ}
  \item \emph{Recall:} say the Tamil for a ticket; a slip of paper, then the Tamil for a map, then say \emph{maṇamakaḷ} again
\end{itemize}

\begin{tcolorbox}[breakable,colback=teal!4,colframe=teal!35!black,title={Before you move on}]
What does \ta{மணமகள்} mean? (\textquotedblleft{}A bride\textquotedblright{}.) Say it once more.
\end{tcolorbox}
\section[\texorpdfstring{\ta{மணமகன்} (maṇamakaṉ)}{மணமகன் (maṇamakaṉ)}]{\ta{மணமகன்} (maṇamakaṉ) — a bridegroom}

\label{lesson:TA-C217-manamakan}



\begin{quote}
Before the new one: say the Tamil for a map, then the Tamil for a bride.
\end{quote}

\subsection*{You'll want to know: \ta{மணமகன்}}
\textbf{\ta{மணமகன்}} — \emph{maṇamakaṉ} — \textquotedblleft{}a bridegroom\textquotedblright{}.

\begin{grammarlens}[title={What you've built}]
One more word for this chapter.
\end{grammarlens}

\subsection*{Your turn}
\begin{itemize}
\raggedright
  \item \emph{Say it:} \emph{maṇamakaṉ}
  \item \emph{Say it:} \emph{maṇamakaṉ}, once more
  \item \emph{Say it:} say \emph{maṇamakaṉ}
  \item \emph{Recall:} say the Tamil for a map, then the Tamil for a bride, then say \emph{maṇamakaṉ} again
\end{itemize}

\begin{tcolorbox}[breakable,colback=teal!4,colframe=teal!35!black,title={Before you move on}]
What does \ta{மணமகன்} mean? (\textquotedblleft{}A bridegroom\textquotedblright{}.) Say it once more.
\end{tcolorbox}
\section[\texorpdfstring{\ta{நிச்சயதார்த்தம்} (niccayatārttam)}{நிச்சயதார்த்தம் (niccayatārttam)}]{\ta{நிச்சயதார்த்தம்} (niccayatārttam) — an engagement (betrothal)}

\label{lesson:TA-C217-niccayatarttam}



\begin{quote}
Before the new one: say the Tamil for a bride, then the Tamil for a bridegroom.
\end{quote}

\subsection*{You'll want to know: \ta{நிச்சயதார்த்தம்}}
\textbf{\ta{நிச்சயதார்த்தம்}} — \emph{niccayatārttam} — \textquotedblleft{}an engagement (betrothal)\textquotedblright{}.

\begin{grammarlens}[title={What you've built}]
One more word for this chapter.
\end{grammarlens}

\subsection*{Your turn}
\begin{itemize}
\raggedright
  \item \emph{Say it:} \emph{niccayatārttam}
  \item \emph{Say it:} \emph{niccayatārttam}, once more
  \item \emph{Say it:} say \emph{niccayatārttam}
  \item \emph{Recall:} say the Tamil for a bride, then the Tamil for a bridegroom, then say \emph{niccayatārttam} again
\end{itemize}

\begin{tcolorbox}[breakable,colback=teal!4,colframe=teal!35!black,title={Before you move on}]
What does \ta{நிச்சயதார்த்தம்} mean? (\textquotedblleft{}An engagement (betrothal)\textquotedblright{}.) Say it once more.
\end{tcolorbox}
\section[\texorpdfstring{\ta{விருந்து} (viruntu)}{விருந்து (viruntu)}]{\ta{விருந்து} (viruntu) — a feast, a party}

\label{lesson:TA-C217-viruntu}



\begin{quote}
Before the new one: say the Tamil for a bridegroom, then the Tamil for an engagement.
\end{quote}

\subsection*{You'll want to know: \ta{விருந்து}}
\textbf{\ta{விருந்து}} — \emph{viruntu} — \textquotedblleft{}a feast, a party\textquotedblright{}.

\begin{grammarlens}[title={What you've built}]
One more word for this chapter.
\end{grammarlens}

\subsection*{Your turn}
\begin{itemize}
\raggedright
  \item \emph{Say it:} \emph{viruntu}
  \item \emph{Say it:} \emph{viruntu}, once more
  \item \emph{Say it:} say \emph{viruntu}
  \item \emph{Recall:} say the Tamil for a bridegroom, then the Tamil for an engagement, then say \emph{viruntu} again
\end{itemize}

\begin{tcolorbox}[breakable,colback=teal!4,colframe=teal!35!black,title={Before you move on}]
What does \ta{விருந்து} mean? (\textquotedblleft{}A feast, a party\textquotedblright{}.) Say it once more.
\end{tcolorbox}
\section[\texorpdfstring{\ta{விழா} (viḻā)}{விழா (viḻā)}]{\ta{விழா} (viḻā) — a festival, a function, a ceremony}

\label{lesson:TA-C217-vila}



\begin{quote}
Before the new one: say the Tamil for an engagement, then the Tamil for a feast.
\end{quote}

\subsection*{You'll want to know: \ta{விழா}}
\textbf{\ta{விழா}} — \emph{viḻā} — \textquotedblleft{}a festival, a function, a ceremony\textquotedblright{}.

\begin{grammarlens}[title={What you've built}]
One more word for this chapter.
\end{grammarlens}

\subsection*{Your turn}
\begin{itemize}
\raggedright
  \item \emph{Say it:} \emph{viḻā}
  \item \emph{Say it:} \emph{viḻā}, once more
  \item \emph{Say it:} say \emph{viḻā}
  \item \emph{Recall:} say the Tamil for an engagement, then the Tamil for a feast, then say \emph{viḻā} again
\end{itemize}

\begin{tcolorbox}[breakable,colback=teal!4,colframe=teal!35!black,title={Before you move on}]
What does \ta{விழா} mean? (\textquotedblleft{}A festival, a function, a ceremony\textquotedblright{}.) Say it once more.
\end{tcolorbox}
